\begin{figure}[ht]
\centering
\begin{tikzpicture}[x=1cm, y=1cm, font=\small]

% Two side-by-side panels: Gaussian sample (points on reference line)
% and heavy-tailed sample (points deviating in both tails).

% ---- Panel A: Gaussian sample ----
\begin{scope}[shift={(0,0)}]
  \draw[->, thick] (-0.2, 0) -- (4.2, 0)
    node[below, font=\scriptsize] {theoretical $z$};
  \draw[->, thick] (0, -0.2) -- (0, 4.2)
    node[left, font=\scriptsize] {observed (sorted)};
  % 45-degree reference line
  \draw[dashed, enggray] (0.2, 0.2) -- (4.0, 4.0);
  % Points approximately on the line
  \foreach \x/\y in {0.4/0.45, 0.8/0.78, 1.1/1.18, 1.5/1.54, 1.8/1.85,
                     2.1/2.05, 2.4/2.45, 2.7/2.68, 3.0/2.95, 3.3/3.30,
                     3.6/3.62} {
    \fill[engblue] (\x, \y) circle (0.06);
  }
  \node[anchor=south, font=\small\bfseries] at (2.0, 4.3)
    {(a) Gaussian sample};
\end{scope}

% ---- Panel B: Heavy-tailed sample ----
\begin{scope}[shift={(5.5,0)}]
  \draw[->, thick] (-0.2, 0) -- (4.2, 0)
    node[below, font=\scriptsize] {theoretical $z$};
  \draw[->, thick] (0, -0.2) -- (0, 4.2)
    node[left, font=\scriptsize] {observed (sorted)};
  % 45-degree reference line
  \draw[dashed, enggray] (0.2, 0.2) -- (4.0, 4.0);
  % Points showing S-shape: low-end values pulled down, high-end up
  \foreach \x/\y in {0.4/0.10, 0.8/0.45, 1.1/0.90, 1.5/1.40, 1.8/1.80,
                     2.1/2.10, 2.4/2.55, 2.7/3.00, 3.0/3.45, 3.3/3.85,
                     3.6/4.10} {
    \fill[engred] (\x, \y) circle (0.06);
  }
  \node[anchor=south, font=\small\bfseries] at (2.0, 4.3)
    {(b) Heavy-tailed sample};
\end{scope}

\end{tikzpicture}
\caption[Q-Q plots against the normal reference]{Q-Q plot of two
samples against a Gaussian reference. The dashed line is the
$y = x$ identity expected when the sample comes from the reference
distribution. Panel (a) shows a $n = 11$ Gaussian sample: the
sorted observations sit close to the line across the full range.
Panel (b) shows a heavy-tailed sample of the same size: the lower
end falls below the line (observed extremes more negative than
Gaussian predicts) and the upper end rises above the line (observed
extremes more positive). The S-shaped deviation diagnoses tails
heavier than Gaussian; a mirror-image S diagnoses lighter tails.}
\label{fig:vol01:ch08:qq-plot}
\end{figure}
